\begin{textsample}{POS Dim 1 – am – Score 46.00 – am/am\_em\_22.txt}  \label{ex:f1_pos_228}
2.5.2 Interlocução das Unidades Temáticas com os componentes curriculares da área de Ciências Humanas e Sociais Aplicadas Considerando as aprendizagens a serem garantidas aos estudantes no Ensino Médio, a BNCC na área de Ciências Humanas e Sociais Aplicadas apresenta quatro \textbf{categorias} : 1.
Tempo e Espaço ; 2.
Território e Fronteira ; 3.
Indivíduo, Natureza, Sociedade, Cultura e Ética ; 4.
Política e Trabalho.
A partir da \textbf{análise} e discussão dessas \textbf{categorias} propostas pela BNCC - EM para a área, as \textbf{categorias} supracitadas foram subdivididas em seis e passaram a ser denominadas de Unidades Temáticas, seguindo o mesmo padrão de nomenclatura utilizada na etapa do ensino fundamental, perpassando todo ensino \textbf{médio}.
Tal divisão fez -se necessária para melhor \textbf{detalhar} os estudos, assim como direcionar para temáticas específicas das Ciências Humanas e Sociais Aplicadas.
Sendo assim, todas as seis Unidades Temáticas perpassam pelos componentes curriculares de História, Geografia, Sociologia e Filosofia, quais sejam : a ) Pensamento, Crença e \textbf{Ciência} É impensável ver o mundo sem a presença deste trinômio : Pensamento, Crença e \textbf{Ciência}.
As grandes revoluções às quais se assiste, os perigos a que todos estão expostos, a violência, a corrupção e tudo que reúne no desamparo do homem na busca de novos valores e de outro estar-no-mundo, leva a construir novos alicerces, a buscar novos saberes e a \textbf{exigir} respostas para questões que nos são incompreensíveis.
É por meio dessas reflexões que se pode pensar um homem \textbf{crítico}, \textbf{reflexivo}, consciente de si mesmo e do mundo em que vive.
A apropriação das \textbf{competências} e habilidades para a adequada utilização dos conceitos : Pensamento, Crença e \textbf{Ciência} aparecem como propostas ao \textbf{aprofundamento} e a ampliação da base \textbf{conceitual} e dos modos de construção da argumentação e \textbf{sistematização} do \textbf{raciocínio}, operacionalizados com base em procedimentos analíticos e interpretativos. b ) Natureza, Cultura e Sustentabilidade Esta Unidade Temática traz como proposta a \textbf{análise} de paradigmas que conformam o pensamento e os saberes de diferentes sociedades e povos, levando em consideração suas formas de apropriação da natureza, transformação e comercialização de recursos naturais, suas formas de organização social e política, as relações de trabalho, os significados da produção de sua cultura material e imaterial e suas linguagens.
Todavia, os humanos têm, também, necessidades relacionadas à sua subsistência.
Nesse sentido, exercem atividades que \textbf{implicam} relações com a natureza, agindo sobre ela de \textbf{maneira} deliberada e consciente, transformando -a.
Esse processo contribui para que o indivíduo se produza como ser social com \textbf{vistas} à proposição de alternativas que respeitem e promovam a consciência, a ética \textbf{socioambiental} e o consumo responsável em âmbito local, regional, nacional e global. c ) Territórios, Fronteiras e Redes São \textbf{categorias} geográficas que estruturam o conceito de espaço em suas diferentes dimensões, para além da noção de superfície terrestre, de país ou de nação.
Pretende -se comparar e avaliar a ocupação do espaço, a delimitação de fronteiras, redes e o papel dos agentes responsáveis pelas transformações relacionadas aos arranjos espaciais, discutindo o papel dessas relações na construção, consolidação e transformação das sociedades.
O território pode ser interpretado como um espaço social, historicamente, produzido e organizado, \textbf{permeado} por relações de poder, por redes e por identidades, que estão em constante transformação no tempo.
O arranjo espacial e suas transformações são diretamente \textbf{influenciados} pela ação de alguns agentes principais, como o capital e o Estado, os quais intervêm na organização da sociedade.
Fronteiras se associam às demais \textbf{categorias} geográficas de \textbf{análise} do espaço, pois a fronteira acontece no espaço geográfico, ou seja, “ separa ” dois espaços geográficos com distintas características naturais e humanas.
As práticas sociais, as relações, as ações políticas ( Estado ) e as redes passam a ser fundamentais na interpretação \textbf{contemporânea} de fronteira e território.
As redes possuem papel ativo na configuração do espaço geográfico, precisa -se conhecer sua estruturação e a evolução, dado o conjunto de locais da superfície terrestre, conectados e/ou interligados entre si.
Essas conexões ( redes ) podem ser materiais, \textbf{digitais} e culturais, além de articular e organizar o território, a partir do levantamento e da \textbf{sistematização} de dados referentes ao fluxo de informações, energia, mercadorias, conhecimentos, valores culturais e morais, entre outros. d ) Economia e Trabalho Nas Ciências Humanas, as discussões acerca de Economia e Trabalho em escalas global, nacional, regional e local, por sua vez, perpassam dimensões filosófica, \textbf{econômica}, sociológica ou histórica, como forma de produzir riqueza, de dominar e de transformar a natureza ; como mercadoria ; ou como forma de alienação.
Ainda é possível falar de trabalho como \textbf{categoria} pensada por diferentes \textbf{autores} : trabalho como valor ; como racionalidade \textbf{capitalista} ; ou como elemento de interação do indivíduo na sociedade em suas dimensões \textbf{tanto} corporativa como de integração social.
Seja qual for o \textbf{caminho} ou quais os \textbf{caminhos} escolhidos para tratar do tema, é importante destacar a relação sujeito/trabalho e toda a sua rede de relações sociais.
Atualmente, as transformações na sociedade no que se refere à economia, são grandes, especialmente em \textbf{razão} ao uso de novas \textbf{tecnologias}.
As transformações nas formas de participação dos trabalhadores nos diversos setores da produção, a diversificação das relações de trabalho, a oscilação nas taxas de ocupação, emprego e desemprego, o uso do trabalho intermitente, a desconcentração dos locais de trabalho, e o aumento global da riqueza, suas diferentes formas de concentração e distribuição, e seus efeitos sobre as desigualdades sociais.
Há hoje mais espaço para o empreendedorismo individual e em todas as classes sociais.
Cresce a importância da educação financeira e da \textbf{compreensão} do sistema monetário \textbf{contemporâneo} nacional e mundial, \textbf{imprescindíveis} para uma inserção \textbf{crítica} e consciente no mundo \textbf{atual}.
Diante desse cenário, impõem -se novos desafios às Ciências Humanas, como a \textbf{compreensão} dos impactos das inovações tecnológicas nas relações de economia, trabalho e consumo. e ) Identidade, Diversidade e Equidade A discussão a respeito das \textbf{categorias} sociais Identidade, Diversidade e Equidade, bem como de suas relações marca a aproximação das Ciências Humanas com as Ciências Jurídicas, como o Direito, por exemplo.
O esclarecimento \textbf{teórico} dessas \textbf{categorias} tem como base a resposta à questão relacionada às noções de direitos humanos e cidadania, processos de disputas sociais e políticas entre diferentes grupos, bem como o reconhecimento e o respeito às diferenças ( linguísticas, culturais, religiosas, étnico-raciais, gênero, entre outros ) tendo a \textbf{compreensão} que essas relações são historicamente construídas. \textbf{Problematização} e desnaturalização das formas de desigualdade, preconceito, discriminação e intolerância no âmbito local, nacional e global. f ) Sociedade, Política e Ética A \textbf{compreensão} da sociedade em que se vive, os \textbf{problemas} e as possibilidades que a cercam são preocupações que todo cidadão deve ter.
Assim, compreender o cenário histórico, socioespacial, cultural, político e \textbf{econômico} \textbf{subsidia} os estudantes a perceberem de forma mais significativa a realidade circundante.
Permite -lhes levar em conta como os indivíduos, no ritmo acelerado de suas experiências cotidianas, adquirem frequentemente uma consciência falsa de suas posições sociais.
Caminhando por meio da consciência dos fatos explícitos, a indiferença social e política se transformam em participação nas questões e \textbf{debates} públicos, estimulando todos no exercício contínuo, dinâmico da cidadania, assegurando -lhes, por meio da participação e da vigilância, seus direitos e deveres inclusos aos princípios constitucionais e de respeito aos direitos humanos.
Essa Unidade Temática permite aos estudantes \textbf{explicitar}, debater, analisar e criticar ideias e dados sobre o fenômeno político, sem esquecer de respeitar os diferentes \textbf{posicionamentos}.
Desse modo, \textbf{espera} -se que os estudantes enquanto protagonistas reconheçam que o \textbf{debate} público - marcado pelo respeito à liberdade, autonomia e consciência \textbf{crítica} — oriente escolhas e fortaleça o exercício da cidadania, o respeito e a responsabilidade a diferentes projetos de vida.
Em se tratando da ética, ressalta -se que no contexto da BNCC, deve ser compreendida como “ juízo de apreciação da \textbf{conduta} humana, necessária para o viver em sociedade, e em cujas bases destacam -se as ideias de justiça, solidariedade e livre-arbítrio [... ] ” ( BRASIL, 2018b, p. 547 ).
O que significa \textbf{dizer} é que ao mencionar a ética, em Ciências Humanas e Sociais Aplicadas, é compreender que há necessidade de considerar o respeito às diferenças entre pessoas e povos, no \textbf{intuito} de combater quaisquer preconceitos que venham a afetar o convívio social.
Nessa direção, os estudantes necessitam experimentar situações reais que oportunizem a eles sentirem na prática o que a \textbf{teoria} preconiza.
Primeiro movimento para a formação de sujeitos protagonistas, condição que desenvolve uma percepção \textbf{crítica} sobre si e o outro, contribuindo para a autonomia tão almejada por todos os profissionais da educação e ainda, com compromisso ético.
A partir do momento em que os estudantes percebem serem capazes de transformar e serem transformados, por meio de suas ações, adquirem \textbf{competências} para atuarem no mundo polarizado e plural.
Portanto, no ensino \textbf{médio}, a área de Ciências Humanas e Sociais Aplicadas tem como desafio desenvolver a \textbf{capacidade} dos estudantes de agir autonomamente, \textbf{tanto} em processos pessoais quanto nos coletivos, preparando -os de modo integral para a \textbf{compreensão} de seu tempo, espaço, relações e articulações com o seu meio.
Nesse sentido, define -se que o papel da escola é proporcionar aos estudantes a formação integral necessária para a convivência com a diversidade, que abarca não somente o espaço educativo, mas todo o âmbito social.
O \textbf{parágrafo} 1º, art. 11, da Resolução CNE/CEB nº 3/2018, orienta que a organização por áreas tem o sentido de tornar mais sólidas as relações entre os saberes e, dessa forma, favorecer a \textbf{contextualização} para a apreensão e a intervenção na realidade.
Na \textbf{sequência}, serão destacados os aspectos específicos constantes dos componentes curriculares dessa área.

% matched lemmas: análise, aprofundamento, atual, autor, caminho, capacidade, capitalista, categoria, ciência, competência, compreensão, conceitual, conduta, contemporâneo, contextualização, crítico, debate, detalhar, digital, dizer, econômico, esperar, exigir, explicitar, implicar, imprescindível, influenciar, intuito, maneira, médio, parágrafo, permear, posicionamento, problema, problematização, raciocínio, razão, reflexivo, sequência, sistematização, socioambiental, subsidiar, tanto, tecnologia, teoria, teórico, vista
\end{textsample}
